% GENERATED FILE. Edit canonical lessons, then run npm run generate:books.
% canonical-source-hash: fnv1a64:945ea01fb99ff05d
% canonical-lessons: SA-C223-yatrikah, SA-C223-margadarsakah, SA-C223-tirtham, SA-C223-asramah, SA-C223-mathah

\chapter{Traveller, Guide, Holy ford, Hermitage, Monastery}
\label{ch:sa-traveller-guide-holy-ford-hermitage-monastery}
\hlchaptermodality{223}

\begin{chapteropening}
\textbf{By the end of this chapter:} \emph{I can say a traveller, a guide, a holy ford, a hermitage and a monastery.}

Complete the last lesson of chapter 223: I can say a traveller, a guide, a holy ford, a hermitage and a monastery.
\end{chapteropening}

\section[\texorpdfstring{\sk{यात्रिकः} (yātrikaḥ)}{यात्रिकः (yātrikaḥ)}]{\sk{यात्रिकः} (yātrikaḥ) — a traveller, a pilgrim}

\label{lesson:SA-C223-yatrikah}



\begin{quote}
Before the new one: say the Sanskrit for a shopkeeper, then the Sanskrit for a newspaper.
\end{quote}

\subsection*{You'll want to know: \sk{यात्रिकः}}
\textbf{\sk{यात्रिकः}} — \emph{yātrikaḥ} — \textquotedblleft{}a traveller, a pilgrim\textquotedblright{}.

\begin{cousinweb}
\emph{yātrika} is built on \emph{yātrā}, \textquotedblleft{}journey, procession\textquotedblright{}, from the root √yā, \textquotedblleft{}to go\textquotedblright{} (MW), with the suffix \emph{-ika}: one on a journey (CDIAL 10457). \emph{Yātrā} lives on as Sindhi \emph{jāṭra}, \textquotedblleft{}body of pilgrims\textquotedblright{} (CDIAL 10456).
\end{cousinweb}

\begin{grammarlens}[title={What you've built}]
One more word for this chapter.
\end{grammarlens}

\subsection*{Your turn}
\begin{itemize}
\raggedright
  \item \emph{Say it:} \emph{yātrikaḥ}
  \item \emph{Say it:} \emph{yātrikaḥ}, once more
  \item \emph{Say it:} say \emph{yātrikaḥ}
  \item \emph{Recall:} say the Sanskrit for a shopkeeper, then the Sanskrit for a newspaper, then say \emph{yātrikaḥ} again
\end{itemize}

\begin{tcolorbox}[breakable,colback=teal!4,colframe=teal!35!black,title={Before you move on}]
What does \sk{यात्रिकः} mean? (\textquotedblleft{}A traveller, a pilgrim\textquotedblright{}.) Say it once more.
\end{tcolorbox}
\section[\texorpdfstring{\sk{मार्गदर्शकः} (mārgadarśakaḥ)}{मार्गदर्शकः (mārgadarśakaḥ)}]{\sk{मार्गदर्शकः} (mārgadarśakaḥ) — a guide}

\label{lesson:SA-C223-margadarsakah}



\begin{quote}
Before the new one: say the Sanskrit for a newspaper, then the Sanskrit for a traveller.
\end{quote}

\subsection*{You'll want to know: \sk{मार्गदर्शकः}}
\textbf{\sk{मार्गदर्शकः}} — \emph{mārgadarśakaḥ} — \textquotedblleft{}a guide\textquotedblright{}.

\begin{cousinweb}
\emph{mārgadarśaka} is MW's \textquotedblleft{}way-shower\textquotedblright{}: \emph{mārga}, \textquotedblleft{}road\textquotedblright{}, with \emph{darśaka}, \textquotedblleft{}one who shows\textquotedblright{}, from the root √dṛś, \textquotedblleft{}to see\textquotedblright{}. \emph{Mārga} first meant the track of a wild animal, from \emph{mṛga}, \textquotedblleft{}deer\textquotedblright{} (CDIAL 10071).
\end{cousinweb}

\begin{grammarlens}[title={What you've built}]
One more word for this chapter.
\end{grammarlens}

\subsection*{Your turn}
\begin{itemize}
\raggedright
  \item \emph{Say it:} \emph{mārgadarśakaḥ}
  \item \emph{Say it:} \emph{mārgadarśakaḥ}, once more
  \item \emph{Say it:} say \emph{mārgadarśakaḥ}
  \item \emph{Recall:} say the Sanskrit for a newspaper, then the Sanskrit for a traveller, then say \emph{mārgadarśakaḥ} again
\end{itemize}

\begin{tcolorbox}[breakable,colback=teal!4,colframe=teal!35!black,title={Before you move on}]
What does \sk{मार्गदर्शकः} mean? (\textquotedblleft{}A guide\textquotedblright{}.) Say it once more.
\end{tcolorbox}
\section[\texorpdfstring{\sk{तीर्थम्} (tīrtham)}{तीर्थम् (tīrtham)}]{\sk{तीर्थम्} (tīrtham) — a holy ford, a place of pilgrimage}

\label{lesson:SA-C223-tirtham}



\begin{quote}
Before the new one: say the Sanskrit for a traveller, then the Sanskrit for a guide.
\end{quote}

\subsection*{You'll want to know: \sk{तीर्थम्}}
\textbf{\sk{तीर्थम्}} — \emph{tīrtham} — \textquotedblleft{}a holy ford, a place of pilgrimage\textquotedblright{}.

\begin{cousinweb}
\emph{tīrtha} is built on the root √tṝ, \textquotedblleft{}to cross over\textquotedblright{} (MW): a ford is where you cross a river (CDIAL 5846). Holy bathing places on rivers became places of pilgrimage. √tṝ is also the root of \emph{vitarati}, \textquotedblleft{}hands out\textquotedblright{}, which you have already met.
\end{cousinweb}

\begin{grammarlens}[title={What you've built}]
One more word for this chapter.
\end{grammarlens}

\subsection*{Your turn}
\begin{itemize}
\raggedright
  \item \emph{Say it:} \emph{tīrtham}
  \item \emph{Say it:} \emph{tīrtham}, once more
  \item \emph{Say it:} say \emph{tīrtham}
  \item \emph{Recall:} say the Sanskrit for a traveller, then the Sanskrit for a guide, then say \emph{tīrtham} again
\end{itemize}

\begin{tcolorbox}[breakable,colback=teal!4,colframe=teal!35!black,title={Before you move on}]
What does \sk{तीर्थम्} mean? (\textquotedblleft{}A holy ford, a place of pilgrimage\textquotedblright{}.) Say it once more.
\end{tcolorbox}
\section[\texorpdfstring{\sk{आश्रमः} (āśramaḥ)}{आश्रमः (āśramaḥ)}]{\sk{आश्रमः} (āśramaḥ) — a hermitage}

\label{lesson:SA-C223-asramah}



\begin{quote}
Before the new one: say the Sanskrit for a guide, then the Sanskrit for a holy ford.
\end{quote}

\subsection*{You'll want to know: \sk{आश्रमः}}
\textbf{\sk{आश्रमः}} — \emph{āśramaḥ} — \textquotedblleft{}a hermitage\textquotedblright{}.

\begin{cousinweb}
\emph{āśrama} is built on \emph{ā-} and the root √śram, \textquotedblleft{}to toil, exert oneself\textquotedblright{} (MW), the root of \emph{śrama}, \textquotedblleft{}toil\textquotedblright{}: a hermitage is a place of religious exertion. It also names the four stages of a life.
\end{cousinweb}

\begin{grammarlens}[title={What you've built}]
One more word for this chapter.
\end{grammarlens}

\subsection*{Your turn}
\begin{itemize}
\raggedright
  \item \emph{Say it:} \emph{āśramaḥ}
  \item \emph{Say it:} \emph{āśramaḥ}, once more
  \item \emph{Say it:} say \emph{āśramaḥ}
  \item \emph{Recall:} say the Sanskrit for a guide, then the Sanskrit for a holy ford, then say \emph{āśramaḥ} again
\end{itemize}

\begin{tcolorbox}[breakable,colback=teal!4,colframe=teal!35!black,title={Before you move on}]
What does \sk{आश्रमः} mean? (\textquotedblleft{}A hermitage\textquotedblright{}.) Say it once more.
\end{tcolorbox}
\section[\texorpdfstring{\sk{मठः} (maṭhaḥ)}{मठः (maṭhaḥ)}]{\sk{मठः} (maṭhaḥ) — a monastery}

\label{lesson:SA-C223-mathah}



\begin{quote}
Before the new one: say the Sanskrit for a holy ford, then the Sanskrit for a hermitage.
\end{quote}

\subsection*{You'll want to know: \sk{मठः}}
\textbf{\sk{मठः}} — \emph{maṭhaḥ} — \textquotedblleft{}a monastery\textquotedblright{}.

\begin{cousinweb}
\emph{maṭha} first meant a hut, the cell of an ascetic or a student (MW), and later a monastery. Its origin is uncertain: CDIAL 9727 sends the reader to \emph{maṇḍapa}, a word it suspects is non-Aryan. It lives on as Hindi \emph{maṛhī}, \textquotedblleft{}hut\textquotedblright{}.
\end{cousinweb}

\begin{grammarlens}[title={What you've built}]
One more word for this chapter.
\end{grammarlens}

\subsection*{Your turn}
\begin{itemize}
\raggedright
  \item \emph{Say it:} \emph{maṭhaḥ}
  \item \emph{Say it:} \emph{maṭhaḥ}, once more
  \item \emph{Say it:} say \emph{maṭhaḥ}
  \item \emph{Recall:} say the Sanskrit for a holy ford, then the Sanskrit for a hermitage, then say \emph{maṭhaḥ} again
\end{itemize}

\begin{tcolorbox}[breakable,colback=teal!4,colframe=teal!35!black,title={Before you move on}]
What does \sk{मठः} mean? (\textquotedblleft{}A monastery\textquotedblright{}.) Say it once more.
\end{tcolorbox}
